\NativeSlot{39}{70}{%
\noindent \textbf{2.}\quad 计算下列各题。（每题3分，共18分）\par\smallskip
\par\noindent%
\ExamTile{.485}{%
% p003_s01

\noindent \((432\div36+45)\times24\)\par


}\hfill%
\ExamTile{.485}{%
% p003_s02

\noindent \(91-(352+272)\div39\)\par


}\par
}
\NativeSlot{172}{57}{%
\par\noindent%
\ExamTile{.485}{%
% p003_s03

\noindent \(592-(118\times4-215)\)\par


}\hfill%
\ExamTile{.485}{%
% p003_s04

\noindent \(76\times[215-(62+78)]\)\par


}\par
}
\NativeSlot{282}{62}{%
\par\noindent%
\ExamTile{.485}{%
% p003_s05

\noindent \([396\div(302-291)]\times58\)\par


}\hfill%
\ExamTile{.485}{%
% p003_s06

\noindent \(184\div[560\div(328-258)]\)\par


}\par
}
\NativeSlot{392}{156}{%
\ExamGroup{四、解决问题。{\normalfont\kaishu\fontsize{10.8bp}{14bp}\selectfont （共32分）}}

% p003_s07

\ExamQuestion{1.}{7}%
四年级18个班在“环保回收·废纸利用”活动中，平均每班回收废纸4千克。学校回收站已回收处理了32千克。剩下的废纸要统一分4批送去再生纸工厂。平均每批送多少千克？先分别从条件、从问题开始分析，填好流程图，再列式解答。\par
\BookFigure{figures/p003_s07-01.tex}
\BookFigure{figures/p003_s07-02.tex}
}
